\documentclass[11pt,a4paper]{article}

\usepackage[utf8]{inputenc}
\usepackage[T1]{fontenc}
\usepackage[turkish]{babel}
\usepackage[a4paper,margin=2.6cm]{geometry}
\usepackage{amsmath}
\usepackage{xcolor}
\usepackage{hyperref}
\usepackage{microtype}
\usepackage{titlesec}
% Belgede hiçbir yerde kalın yazı kullanılmıyor; bölüm başlıkları da
% varsayılan kalın biçimden çıkarılıyor.
\titleformat{\section}{\normalfont\Large\color{ink}}{}{0pt}{}
\titlespacing*{\section}{0pt}{22pt plus 5pt}{9pt}

\definecolor{myco}{HTML}{0E8F76}
\definecolor{ember}{HTML}{C1670E}
\definecolor{ink}{HTML}{1E2622}
\definecolor{muted}{HTML}{727C75}

\hypersetup{colorlinks=true, linkcolor=myco, urlcolor=myco,
  pdftitle={Mycellium-Aegis - Yapay Zeka Gerekcesi},
  pdfauthor={Mycellium-Aegis}}

% Soru-cevap bicimi: kalin yazi kullanilmiyor, ayrim kucuk kapital
% etiketler ve soru satirindaki italik ile saglaniyor.
\newcommand{\soru}[1]{%
  \par\vskip 15pt plus 4pt minus 2pt
  \noindent\textcolor{ember}{\textsc{Soru}}\quad\emph{#1}%
  \par\vskip 6pt\noindent\textcolor{myco}{\textsc{Cevap}}\quad\ignorespaces}

\setlength{\parskip}{6pt plus 2pt minus 1pt}
\setlength{\parindent}{0pt}
\linespread{1.06}

\title{\normalfont\LARGE Mycellium-Aegis\\[6pt]
  \large Yapay Zeka Neden Var, Nasıl Çalışıyor, Nerede Zorlanıyor\\[5pt]
  \normalsize\textcolor{muted}{İş fikri değerlendirmesi için soru-cevap notu}}
\author{\normalfont\normalsize Mycellium-Aegis ekibi}
\date{\normalfont\small TÜBİTAK 1812 BİGG \quad·\quad Temmuz 2026}

\begin{document}
\maketitle
\thispagestyle{empty}

\textcolor{muted}{Bu not, sistemin yapay zeka tarafını teknik olmayan bir
değerlendiriciye anlatmak için hazırlandı. Matematiksel türetmeler ve
donanım ayrıntıları ayrı belgelerde duruyor; burada yalnızca neyi neden
yaptığımız, hangi kararı hangi gerekçeyle verdiğimiz ve nerede
zorlandığımız anlatılıyor. Ölçülmüş olanla henüz hedef olanı birbirinden
ayırmaya özellikle dikkat ettik.}

\vskip 10pt

\section*{Sorunun kendisi}

\soru{Toprağa bir sensör gömüp bir eşik değeri koymak varken neden yapay
zekaya ihtiyaç duyuyorsunuz? Bu bir aşırı mühendislik değil mi?}
Bu itiraz haklı bir yerden geliyor, o yüzden önce neyi ölçtüğümüzü
söylemek gerekiyor. Ölçtüğümüz şey bir yangın değil, bir canlının
strese verdiği tepki. Miselyum ağı üzerinden okuduğumuz gerilim, ısı
şoku dışında pek çok şeye de tepki veriyor: yağmur sonrası nem sıçraması,
toprağın gün içinde ısınıp soğuması, yakınından geçen bir hayvan,
kökün büyümesi, elektrotun etrafındaki toprağın oturması. Bunların
hepsi aynı kanalda, aynı büyüklük mertebesinde görünüyor. Tek bir
gerilim eşiği koyduğunuzda ya eşiği yükseltip gerçek yangını kaçırırsınız
ya da alçaltıp sistemi haftada birkaç kez boşuna alarm veren, kimsenin
ciddiye almadığı bir kutuya dönüştürürsünüz. Ormancılık tarafında bir
sistemin öldüğü nokta tam olarak burasıdır; teknik başarısızlıktan çok,
güvenilirliğini kaybetmesinden ölür.

Yapay zekayı bu yüzden kullanıyoruz. İşi, sinyali bir eşikle kesmek
yerine, sinyalin zaman içindeki biçimine bakıp ``bu ısı şokuna benziyor
mu, yoksa nem sıçramasına mı'' sorusunu cevaplamak. Yani yapay zeka
burada bir süs değil, ürünü kullanılabilir kılan bileşen. Onsuz da bir
prototip yapılabilir, ama sahada bir yıl ayakta kalacak bir ürün
yapılamaz.

\soru{Peki sistem uçtan uca nasıl çalışıyor? Bir yangın başladığında
sırasıyla ne oluyor?}
Yerin on sekiz santimetre altındaki biyoelektrot, miselyum ağının
gerilimini sürekli okuyor. Bu ham okuma çok küçük ve çok gürültülü
olduğu için önce düğümün üzerinde temizleniyor: uzun vadeli sürüklenme
çıkarılıyor, yüksek frekanslı elektriksel parazit süzülüyor. Ardından
sinyalden birkaç sayı çıkarılıyor. Bunlar sinyalin baskın frekansı,
genliğinin ne kadar düştüğü, aynı kablodaki sıcaklık probunun ne hızla
ısındığı ve elektrot direncinin ne yönde değiştiği. Yani ham dalga
biçimi değil, onun birkaç özeti taşınıyor.

Bu özetler düğümün belleğinde bir süre birikiyor ve bir zaman serisi
modeline veriliyor. Model, son yarım saatlik davranışa bakıp bir
olasılık üretiyor. Bu olasılık tek başına alarm vermiyor; yanına
fiziksel bir koşul listesi konuyor. Sıcaklık yeterince hızlı yükselmiş
olmalı, toprak direnci kuruma yönünde değişmiş olmalı ve biyoelektrik
imza kendi bandında görünmüş olmalı. Üçü birden sağlanmadan alarm
çıkmıyor. Alarm çıktığında LoRa üzerinden yüzey düğümüne, oradan ağ
geçidine ve buluta gidiyor; panoda konumu ve hangi koşulların
tetiklendiği görünüyor.

\soru{Modeli neden düğümün üzerinde çalıştırıyorsunuz? Buluta gönderip
orada işlemek daha kolay olmaz mıydı?}
Kolay olurdu ama sistem sahada yaşayamazdı. Yeraltı düğümü bir pille
çalışıyor ve enerji bütçesinin neredeyse tamamı telsiz iletiminde
harcanıyor; hesaplama yanında çok küçük kalıyor. Ham sinyali sürekli
yukarı göndermek, pili haftalar yerine günler mertebesine indirir.
İkinci sebep kapsama: orman içinde bağlantı sürekli değil, düğüm
saatlerce yalnız kalabiliyor. Karar buluta bağlıysa, bağlantının
koptuğu anda sistem kör oluyor. Bu yüzden ağır iş aşağıda yapılıyor,
yukarıya yalnızca birkaç sayı ve bir karar çıkıyor. Bulut tarafı da
işlevsiz değil; modelin yeni sürümleri orada eğitiliyor ve düğümlere
indiriliyor, ayrıca birden fazla düğümün aynı anda ne gördüğü orada
birleştiriliyor.

\section*{Tasarım kararları ve gerekçeleri}

\soru{Neden bu tip bir zaman serisi modeli seçtiniz? Daha güçlü ve
güncel modeller varken bu tercih geri kalmış görünmüyor mu?}
Model seçimini üç kısıt belirledi. Birincisi, elimizdeki veri az.
Etiketli yangın verisi doğası gereği kıt, çünkü kontrollü koşulda
yangın üretmek pahalı ve sınırlı. Az veriyle çalışırken büyük modeller
ezberliyor, küçük modeller genelliyor. İkincisi, model bir mikrodenetleyici
üzerinde, kilobaytlar mertebesinde bellekle ve pil bütçesiyle çalışacak.
Üçüncüsü, kararın neden verildiğini açıklayabilmemiz gerekiyor; orman
idaresine ``model öyle dedi'' demek yeterli bir gerekçe değil. Bu üç
kısıt bizi mütevazı boyutlu, sırayı öğrenen bir modele götürdü. Amacımız
literatürdeki en güçlü modeli kullanmak değil, sahada bir yıl boyunca
aynı davranışı gösterecek modeli kullanmak.

Şunu da açıkça söyleyelim: bu model henüz eğitilmedi. Elimizde şu an
laboratuvarda ölçülmüş sinyal davranışı ve fiziksel karar kuralı var.
Modelin eğitimi ve doğrulanması projenin altıncı ile dokuzuncu ayları
arasındaki iş paketi. Bugün sunduğumuz şey eğitilmiş bir model değil,
eğitilebileceğini gösteren ölçüm ve bunun üzerine kurulmuş bir plan.

\soru{Modelin kararını neden fiziksel bir koşul listesiyle
sınırlıyorsunuz? Modele güvenmiyor musunuz?}
Güven meselesi değil, sorumluluk meselesi. Bir yangın alarmı, insan ve
araç hareket ettiren bir sinyal. Böyle bir kararı tek bir öğrenilmiş
modele bırakmak, modelin daha önce görmediği bir durumla karşılaştığında
ne yapacağını bilmemek demek. Fiziksel koşullar burada bir güvenlik ağı
işlevi görüyor: sıcaklık hiç yükselmediyse, toprak direnci hiç
değişmediyse, model ne derse desin alarm çıkmıyor. Bu, yanlış alarmı
bastırmanın en ucuz ve en denetlenebilir yolu. Aynı zamanda kararı
açıklanabilir kılıyor, çünkü alarm ekranında hangi koşulun ne zaman
tetiklendiği tek tek görünüyor.

Bunun bir bedeli var ve onu da söylemek gerekir: koşulların hepsini
birden aramak, duyarlılığı bir miktar düşürüyor. Çok yavaş gelişen,
sıcaklığı yeterince hızlı yükseltmeyen bir için için yanmayı
kaçırabiliriz. Bu bilinçli bir tercih. Erken aşamada kaçırılan bir
olayın maliyeti, güvenilirliğini kaybetmiş bir sistemin maliyetinden
düşük.

\soru{Yanlış alarm oranı için verdiğiniz hedef neye dayanıyor?}
Yüzde birin altı hedefi, ölçülmüş bir sonuç değil, bir tasarım hedefi
ve bunu böyle sunuyoruz. Dayanağı şu: bir ağ geçidine yüz düğüm bağlı
ve her düğüm günde birkaç kez karar üretiyorsa, yüzde birlik bir hata
bile ekipleri haftada birkaç kez boşuna yola çıkarır. Operasyonu
sürdürülebilir kılan eşik burada. Hedefe fiziksel koşulların birlikte
aranmasıyla ve modelin buna göre ayarlanmasıyla ulaşmayı planlıyoruz;
gerçek oran ancak sahada üç aylık kesintisiz veriyle ölçülebilir, o da
dokuzuncu ile on ikinci aylar arasındaki iş paketi.

\section*{Zorluklar ve bunlara yanıtımız}

\soru{Toprak tipi, mevsim, mantar türü ve nem koşulları her sahada
farklı. Bir sahada eğitilmiş model başka sahada çalışır mı?}
Doğrudan çalışmaz, bunu baştan kabul ediyoruz. Bu yüzden modeli mutlak
değerler üzerine değil, değişim üzerine kuruyoruz. Bir sahanın
elektrot direncinin kaç kilo ohm olduğu sahadan sahaya değişir; ama o
direncin bir saat içinde hangi yönde ve hangi hızda değiştiği çok daha
taşınabilir bir bilgidir. Aynı şekilde gerilimin mutlak seviyesi değil,
dinlenim durumuna göre ne kadar saptığı anlamlı. Kurulumdan sonra her
düğüm birkaç gün boyunca kendi normalini öğreniyor ve kararlarını bu
kişisel taban çizgisine göre veriyor. Bu, modelin sahaya taşınmasını
kolaylaştıran temel tasarım kararı.

Yine de bunun sınırı var. Çok farklı bir iklim kuşağında, çok farklı
bir tür kompozisyonunda yeniden kalibrasyon gerekecektir. Ticari planda
bunu bir zayıflık değil, bir hizmet kalemi olarak konumlandırıyoruz:
kurulum, kalibrasyon ve ilk sezon takibi paketin parçası.

\soru{Laboratuvarda ölçtüğünüz sinyalin gerçekten ısı kaynaklı olduğunu
nereden biliyorsunuz? Ölçüm hatası olamaz mı?}
Olabilir ve bu ihtimali kapatmak için birkaç şey yapıyoruz. Deneyde
gördüğümüz şey, dinlenim halinde belirli bir frekansta salınan sinyalin,
ısı şoku uygulandığında hem yavaşlaması hem de genliğinin belirgin bir
alt sınıra inmesiydi. Bu iki değişimin birlikte olması, tek başına bir
kayma olmasından daha güçlü bir kanıt. Bununla birlikte ölçüm
düzeneğimizde bir zayıflık tespit ettik ve bunu saklamıyoruz: örnekleme
hızımız, ölçtüğümüz frekansa fazla yakındı. Bu koşulda gözlenen
frekansın gerçek mi yoksa örnekleme kaynaklı bir yanılsama mı olduğu
kesin ayrılamaz. Deneyi daha yüksek örnekleme hızıyla ve girişe
takma-ad önleyici bir süzgeç koyarak tekrarlamayı iş planına aldık.
Bir ölçümü savunmanın yolu onu tekrarlamaktır; ilk işimiz bu.

\soru{Sistemin en zayıf halkası nedir?}
Bugün itibarıyla en zayıf halka model değil, verinin kendisi. Etiketli
yangın verisi az ve üretmesi pahalı. Buna karşı üç koldan çalışıyoruz.
İklim kabininde kontrollü ısı şoku uygulayarak etiketli veri üretiyoruz.
Sahaya yerleştirilen düğümlerden yangın olmayan uzun süreli veri
topluyoruz, çünkü yanlış alarmı bastırmak için asıl ihtiyaç duyulan şey
normalin ne olduğunu iyi bilmek. Üçüncü olarak, fiziksel çözücümüzle
sentetik senaryolar üretip modeli bunlarla ön eğitime tabi tutuyoruz.
Sentetik veri gerçeğin yerini tutmaz, ama modelin sıfırdan başlamamasını
sağlar.

İkinci zayıf halka, sensörün derinliği konusunda belgelerimizle
hesabımız arasında çıkan bir tutarsızlıktı. Isı iletim çözümünü
çalıştırdığımızda, sıcaklık probunun miselyumla aynı derinlikte
durmasının erken uyarıyı imkânsız kıldığını gördük; ısı cephesi o
derinliğe tutuşmadan sonra varıyor. Çözüm basit oldu: biyoelektrot
miselyumun yaşadığı derinlikte kalıyor, sıcaklık probu aynı kablo
üzerinde yüzeye yakın bir noktaya alınıyor. Bu değişiklik maliyeti üç
doların altında artırıyor ve erken uyarı penceresini açıyor. Bunu
burada anmamızın sebebi şu: modelimiz bir hatayı yakaladı, biz de
tasarımı ona göre düzelttik. Sayısal ikizi bunun için kurduk.

\section*{Ticari taraf}

\soru{Rakipleriniz de yapay zeka kullandığını söylüyor. Sizinki neden
farklı?}
Fark modelde değil, modele giren veride. Kamera tabanlı sistemler
görüntüden duman ayırt etmeye çalışıyor; gaz sensörleri havadaki
partikülü sınıflandırıyor. İkisi de yangının çıktısını işliyor. Biz
yangının nedenini, yani toprakta yayılan ısı cephesini ve canlının ona
verdiği tepkiyi işliyoruz. Aynı modelleme tekniğini kullansanız bile,
daha erken bir sinyale bakıyor olmak yapısal bir avantaj. Görüş hattı,
rüzgâr yönü ve hava koşulu bizim için bir kısıt değil, çünkü ölçüm
yerin altında yapılıyor.

\soru{Bu yapay zeka savunulabilir bir varlık mı, yoksa kopyalanabilir mi?}
Modelin mimarisi kopyalanabilir; bunda özel bir şey yok. Kopyalanamayan
şey, o modeli işe yarar kılan iki unsur. Birincisi biyomimetik elektrot
ve onun toprakla kurduğu düşük empedanslı arayüz; iyi sinyal
alamıyorsanız hangi modeli kullandığınızın önemi kalmıyor. Bu tarafta
patent başvurusu planlıyoruz. İkincisi, sahadan biriken normal davranış
verisi. Kurulu her düğüm, zamanla o bölgenin normalini daha iyi bilen
bir veri kümesi üretiyor ve bu küme bizde kalıyor. Rekabet avantajı
zamanla modelde değil, veride birikiyor.

\soru{On sekiz ayın sonunda yapay zeka tarafında elinizde somut olarak
ne olacak?}
Üç şey. Kontrollü koşulda üretilmiş, etiketli bir ısı şoku veri kümesi
ve bunun yanında sahadan toplanmış üç aylık kesintisiz normal davranış
verisi. Bu veriyle eğitilmiş, mikrodenetleyici üzerinde çalışan ve
enerji bütçesine sığdığı ölçülmüş bir karar modeli. Ve bu modelin saha
koşulunda ürettiği yanlış alarm oranının, tahmin değil ölçüm olarak
raporlanması. Üçüncüsü en önemlisi, çünkü satın alma kararını verecek
kurumun soracağı ilk soru bu olacak.

\vskip 18pt
\hrule
\vskip 8pt
\textcolor{muted}{\small Bu belgede ölçülmüş olarak sunulan tek veri
kümesi, laboratuvarda \emph{Pleurotus ostreatus} ağı üzerinde yapılan
ısı şoku deneyidir. Yanlış alarm oranı, batarya ömrü ve saha
başarımına ilişkin sayılar hedef veya modeldir; ilgili iş paketleri
tamamlandığında ölçümle değiştirilecektir.}

\end{document}
